\documentclass[12pt]{article}

\usepackage{amsmath, amssymb, amsthm}
\usepackage{geometry}
\usepackage{booktabs}
\usepackage{hyperref}
\usepackage{natbib}
\usepackage{microtype}
\usepackage{bm}

\geometry{margin=1.2in}

\newtheorem{definition}{Definition}
\newtheorem{proposition}{Proposition}
\newtheorem{remark}{Remark}

\title{Residual Analysis for the Adjusted Score-Driven Factor Model:\\
The Bucketed Case}
\date{}

\begin{document}

\maketitle

\begin{abstract}
We develop a systematic residual analysis for the adjusted score-driven factor model
of Zou, Lin, and Lucas (2025) applied to the 24-dimensional bucketed implied volatility
surface. The dimension is fixed at $N = 24$ throughout the sample, which allows the full
battery of univariate, cross-sectional, and multivariate panel tests to be applied without
the complications arising from a time-varying grid. We cover standardization, marginal
distributional tests, serial dependence tests, cross-sectional dependence tests, the multivariate
portmanteau test, and ARCH-LM tests. Throughout, null hypotheses are stated precisely and
the asymptotic distributions of each statistic are derived or referenced.
\end{abstract}

\tableofcontents
\newpage

%------------------------------------------------------------------
\section{Setup and Notation}
%------------------------------------------------------------------

We work with the bucketed dataset in which the log implied volatility surface is represented
at each date $t = 1, \ldots, T$ by a fixed $N = 24$ dimensional vector
\begin{equation}
    y_t = \log IV_t \in \mathbb{R}^{24},
\end{equation}
where each entry corresponds to a unique (moneyness, time-to-maturity) bucket. The
measurement equation of the adjusted score-driven factor model is
\begin{equation}
    y_t = M_t \beta_t + \varepsilon_t, \qquad
    \varepsilon_t \sim h\!\left(\varepsilon_t \mid \Sigma_t;\, \nu\right),
    \label{eq:measurement}
\end{equation}
where $M_t \in \mathbb{R}^{N \times p}$ collects the factor loadings evaluated at the bucket
representative contracts, $\beta_t \in \mathbb{R}^p$ is the vector of time-varying factors, and
the measurement noise covariance is
\begin{equation}
    \Sigma_t = H_t + M_t C M_t^\top,
    \label{eq:sigma}
\end{equation}
with $H_t$ a diagonal matrix of idiosyncratic variances and $C \in \mathbb{R}^{p \times p}$ the
additional static parameter matrix introduced to align the predictive covariance of the
score-driven model with that of its state-space counterpart. In the bucketed case $M_t$
is effectively time-invariant because each bucket is always represented by its midpoint
contract, so we write $M_t \equiv M$ and $\Sigma_t \equiv \Sigma$ for the purpose of
the residual analysis. The state transition is
\begin{equation}
    \beta_{t+1} = (I_p - B)\bar{\beta} + B\beta_t + \xi_t,
\end{equation}
where $\xi_t$ is the score-driven increment computed from data up to time $t$.

After estimating the static parameter vector $\psi$ by maximum likelihood and obtaining
the filtered factor path $\{\hat{\beta}_t\}_{t=1}^T$, the raw residuals are
\begin{equation}
    \hat{\varepsilon}_t = y_t - M\hat{\beta}_t \in \mathbb{R}^{24}.
\end{equation}
The goal of the residual analysis is to determine whether $\{\hat{\varepsilon}_t\}$ is
consistent with the maintained distributional assumption $h(\,\cdot \mid \hat{\Sigma};\,\hat{\nu})$,
and in particular whether the residuals are approximately i.i.d.\ after standardization.

%------------------------------------------------------------------
\section{Standardized Residuals}
%------------------------------------------------------------------

\subsection{Individual Standardization}

The marginal variance of the $i$-th residual implied by \eqref{eq:sigma} is
$[\Sigma]_{ii}$. Define the individually standardized residuals
\begin{equation}
    z_{i,t} = \frac{\hat{\varepsilon}_{i,t}}{\sqrt{[\hat{\Sigma}]_{ii}}},
    \qquad i = 1,\ldots,24,\quad t = 1,\ldots,T.
    \label{eq:zit}
\end{equation}
Under correct specification and a Gaussian assumption, $z_{i,t} \overset{iid}{\sim} \mathcal{N}(0,1)$.
Under Student's $t$ with $\nu$ degrees of freedom, the standardized residuals follow a
unit-variance $t$ distribution, i.e.\ $z_{i,t} \overset{iid}{\sim} t\!\left(0, 1, \hat{\nu}\right)$.
Individual standardization is sufficient for univariate tests but not for multivariate
tests, which require full Mahalanobis standardization as described below.

\subsection{Mahalanobis Standardization}

For multivariate tests one works with the Mahalanobis distance
\begin{equation}
    d_t^2 = \hat{\varepsilon}_t^\top \hat{\Sigma}^{-1} \hat{\varepsilon}_t.
    \label{eq:mahal}
\end{equation}
Under correct Gaussian specification, $d_t^2 \overset{iid}{\sim} \chi^2(N)$ with $N = 24$.
It is convenient to standardize further as
\begin{equation}
    w_t = \frac{d_t^2 - N}{\sqrt{2N}},
\end{equation}
so that $w_t \overset{iid}{\approx} \mathcal{N}(0,1)$ for large $N$ by the central limit theorem.

For the Student's $t$ case, the Mahalanobis distance has the scaled $F$ distribution
\begin{equation}
    \frac{d_t^2 / N}{(\hat{\nu} - 2)/\hat{\nu}} \sim F(N, \hat{\nu}),
\end{equation}
and the same $\chi^2$ approximation applies as $\hat{\nu} \to \infty$.

\begin{remark}
    In the bucketed case $\hat{\Sigma}$ is $24 \times 24$ and time-invariant after
    plugging in estimated parameters, so inverting it once suffices. In the non-bucketed
    case $\Sigma_t$ changes dimension with $t$, making \eqref{eq:mahal} substantially
    harder to use.
\end{remark}

\subsection{Probability Integral Transform}

For the Gaussian case, define the probability integral transform (PIT) values
\begin{equation}
    u_{i,t} = \Phi(z_{i,t}),
    \label{eq:pit}
\end{equation}
where $\Phi$ is the standard normal CDF. Under correct specification,
$u_{i,t} \overset{iid}{\sim} \mathrm{Uniform}(0,1)$. For the Student's $t$ case, replace
$\Phi$ with the CDF of the unit-variance $t(\hat{\nu})$ distribution. PIT values are the
natural input for distribution tests described in Section~\ref{sec:marginal}.

%------------------------------------------------------------------
\section{Marginal Distribution Tests}
\label{sec:marginal}
%------------------------------------------------------------------

\subsection{Normality Tests on $z_{i,t}$}

\paragraph{Jarque-Bera test.}
For each series $i$, compute the sample skewness $\hat{s}_i$ and excess kurtosis $\hat{k}_i$
of $\{z_{i,t}\}_{t=1}^T$. The Jarque-Bera statistic is
\begin{equation}
    JB_i = \frac{T}{6}\hat{s}_i^2 + \frac{T}{24}\hat{k}_i^2
    \overset{d}{\to} \chi^2(2) \quad \text{under } H_0.
\end{equation}
Run this for each $i = 1,\ldots,24$. To obtain a joint test across all series, apply
Fisher's combination:
\begin{equation}
    \mathcal{F}_{JB} = -2\sum_{i=1}^{24} \log p_i^{JB}
    \overset{d}{\to} \chi^2(48),
\end{equation}
where $p_i^{JB}$ is the p-value of $JB_i$.

\paragraph{Kolmogorov-Smirnov test.}
For each $i$, let $\hat{F}_i$ denote the empirical CDF of $\{z_{i,t}\}$. The KS statistic is
\begin{equation}
    KS_i = \sup_{x \in \mathbb{R}} \left|\hat{F}_i(x) - \Phi(x)\right|.
\end{equation}
The distribution of $\sqrt{T} \cdot KS_i$ under $H_0$ is the Kolmogorov distribution and
does not depend on the null parameters. Combine across $i$ via Fisher's method as above.

\paragraph{Berkowitz test on PIT values.}
Transform the PIT values \eqref{eq:pit} back via $r_{i,t} = \Phi^{-1}(u_{i,t}) = z_{i,t}$.
If the $u_{i,t}$ are i.i.d.\ Uniform$(0,1)$, the $r_{i,t}$ should be i.i.d.\ $\mathcal{N}(0,1)$.
Berkowitz (2001) proposes fitting an AR(1) to $r_{i,t}$:
\begin{equation}
    r_{i,t} = \mu_i + \rho_i r_{i,t-1} + \eta_{i,t}, \qquad \eta_{i,t} \sim \mathcal{N}(0, \sigma_i^2),
\end{equation}
and testing $H_0: \mu_i = 0,\, \rho_i = 0,\, \sigma_i^2 = 1$ jointly via a likelihood ratio
statistic:
\begin{equation}
    LR_i^{Berk} = -2\left[\ell_0(r_{i,\cdot}) - \ell_1(\hat{\mu}_i, \hat{\rho}_i, \hat{\sigma}_i^2; r_{i,\cdot})\right]
    \overset{d}{\to} \chi^2(3),
\end{equation}
where $\ell_0$ is the log-likelihood under $H_0$ and $\ell_1$ is the unrestricted
log-likelihood. Apply Fisher's method to combine across $i$.

%------------------------------------------------------------------
\section{Serial Dependence Tests}
%------------------------------------------------------------------

\subsection{Ljung-Box Test per Series}

For each series $i$, compute the sample autocorrelation of $z_{i,t}$ at lag $\ell$:
\begin{equation}
    \hat{\rho}_{i,\ell} = \frac{\sum_{t=\ell+1}^T z_{i,t} z_{i,t-\ell}}
    {\sum_{t=1}^T z_{i,t}^2}.
\end{equation}
The Ljung-Box statistic at maximum lag $L$ is
\begin{equation}
    LB_i(L) = T(T+2) \sum_{\ell=1}^{L} \frac{\hat{\rho}_{i,\ell}^2}{T - \ell}
    \overset{d}{\to} \chi^2(L) \quad \text{under } H_0.
    \label{eq:lb}
\end{equation}
Repeat for $z_{i,t}^2$ to detect ARCH-type dependence in higher moments, replacing
$z_{i,t}$ with $z_{i,t}^2 - 1$ in the autocorrelation formula.

\subsection{Combined Serial Dependence Test Across Series}

Denote by $p_i^{LB}$ the p-value of \eqref{eq:lb} for series $i$. Fisher's combination
statistic is
\begin{equation}
    \mathcal{F}_{LB} = -2 \sum_{i=1}^{24} \log p_i^{LB}
    \overset{d}{\to} \chi^2(48) \quad \text{under } H_0.
    \label{eq:fisher_lb}
\end{equation}
Apply the same procedure separately to $z_{i,t}^2$ to obtain $\mathcal{F}_{LB^2}$. The
two statistics together test the null of no serial dependence in the first two moments of the
standardized residuals, jointly across all 24 series.

\begin{remark}
    Fisher's method assumes independence across the $p_i^{LB}$, which is not guaranteed
    when the series share common factors. A conservative correction is to apply the
    Bonferroni-Holm step-down procedure: sort the 24 p-values in ascending order
    $p_{(1)} \leq \cdots \leq p_{(24)}$ and reject the null for series $j$ if
    $p_{(j)} \leq \alpha / (25 - j)$. The family-wise error rate is controlled at level
    $\alpha$.
\end{remark}

%------------------------------------------------------------------
\section{Cross-Sectional Dependence Tests}
%------------------------------------------------------------------

\subsection{Pesaran CD Test}

Under the null of no cross-sectional correlation, the off-diagonal elements of the
long-run cross-sectional covariance of $\{z_t\}$ are zero. Define the pairwise sample
correlation between series $i$ and $j$:
\begin{equation}
    \hat{\rho}_{ij} = \frac{\sum_{t=1}^T z_{i,t} z_{j,t}}
    {\sqrt{\sum_{t=1}^T z_{i,t}^2 \cdot \sum_{t=1}^T z_{j,t}^2}}.
\end{equation}
Pesaran's (2004) CD statistic is
\begin{equation}
    CD = \sqrt{\frac{2T}{N(N-1)}} \sum_{i=1}^{N-1} \sum_{j=i+1}^{N} \hat{\rho}_{ij}
    \overset{d}{\to} \mathcal{N}(0,1) \quad \text{under } H_0,
    \label{eq:cd}
\end{equation}
where $N = 24$. The null hypothesis is that all pairwise contemporaneous correlations
are zero. Under the alternative that the factor model has failed to absorb all
cross-sectional dependence, $\hat{\rho}_{ij} > 0$ on average and $CD$ diverges to
$+\infty$.

\begin{remark}
    The CD test has power against both omitted factors and misspecified factor loadings.
    A significant CD test after estimating the factor model is therefore informative: it
    signals that the adjusted covariance $\hat{\Sigma} = \hat{H} + M\hat{C}M^\top$ has not
    fully captured the cross-sectional dependence. In practice, one expects residual
    dependence to vanish after the covariance adjustment is applied.
\end{remark}

\subsection{Testing the Full Cross-Sectional Covariance Matrix}

A stronger test examines whether the full sample cross-sectional covariance
\begin{equation}
    \hat{\Gamma}_0 = \frac{1}{T} \sum_{t=1}^T z_t z_t^\top
\end{equation}
equals the identity matrix $I_{24}$. Define the statistic
\begin{equation}
    S_{cov} = \frac{T}{2} \left[\mathrm{tr}\!\left(\hat{\Gamma}_0^2\right) - N - \frac{1}{T}
    \left(\mathrm{tr}(\hat{\Gamma}_0)\right)^2\right]
    \overset{d}{\to} \chi^2\!\left(\frac{N(N-1)}{2}\right)
    \quad \text{under } H_0,
    \label{eq:covtest}
\end{equation}
which tests $H_0: \Gamma_0 = I_N$ against the general alternative of any non-identity
covariance structure. With $N = 24$ the degrees of freedom are $24 \times 23 / 2 = 276$.

%------------------------------------------------------------------
\section{Multivariate Portmanteau Test}
%------------------------------------------------------------------

\subsection{Hosking's Statistic}

The multivariate analogue of the Ljung-Box test is due to Hosking (1980). Define the
lag-$\ell$ cross-covariance matrix of the standardized residuals:
\begin{equation}
    \hat{\Gamma}(\ell) = \frac{1}{T} \sum_{t=\ell+1}^T z_t z_{t-\ell}^\top,
    \qquad \ell = 1, \ldots, L.
\end{equation}
Hosking's multivariate portmanteau statistic is
\begin{equation}
    Q_M(L) = T^2 \sum_{\ell=1}^{L} \frac{1}{T-\ell}
    \mathrm{tr}\!\left[\hat{\Gamma}(\ell)^\top \hat{\Gamma}(0)^{-1}
    \hat{\Gamma}(\ell) \hat{\Gamma}(0)^{-1}\right]
    \overset{d}{\to} \chi^2(N^2 L) \quad \text{under } H_0.
    \label{eq:hosking}
\end{equation}
With $N = 24$ and $L = 10$ (a common choice), the degrees of freedom are
$576 \times 10 = 5760$, which can make the test conservative in finite samples. A
practical alternative is the Li-McLeod correction:
\begin{equation}
    Q_{LM}(L) = Q_M(L) + \frac{N^2 L(L+1)}{2T}
    \overset{d}{\to} \chi^2(N^2 L),
\end{equation}
which has better size in finite samples than Hosking's original formulation.

\begin{remark}
    The test in \eqref{eq:hosking} uses $\hat{\Gamma}(0)$ rather than $I_N$. If
    $\hat{\Gamma}(0) \neq I_N$ it is advisable to first test cross-sectional
    dependence via \eqref{eq:covtest}, and then run \eqref{eq:hosking}
    on the fully standardized residuals $\tilde{z}_t = \hat{\Gamma}(0)^{-1/2} z_t$,
    which by construction satisfy $T^{-1}\sum_t \tilde{z}_t \tilde{z}_t^\top = I_N$.
\end{remark}

\subsection{Recommended Lag Length}

A common rule is $L = \lfloor T^{1/3} \rfloor$. For $T = 500$ (the rolling window used
in the empirical application), this gives $L \approx 8$. The test should be run for
several values of $L$ to ensure robustness, as the power profile against different
alternatives depends on the lag length chosen.

%------------------------------------------------------------------
\section{ARCH-LM Test}
%------------------------------------------------------------------

\subsection{Univariate ARCH-LM per Series}

For each series $i$, the Engle (1982) LM test for ARCH effects of order $q$ regresses
$z_{i,t}^2$ on its own $q$ lags:
\begin{equation}
    z_{i,t}^2 = \alpha_{i,0} + \sum_{j=1}^{q} \alpha_{i,j} z_{i,t-j}^2 + \eta_{i,t}.
    \label{eq:archreg}
\end{equation}
The LM statistic is $ARCH_i(q) = T \cdot R_i^2$, where $R_i^2$ is the coefficient of
determination of \eqref{eq:archreg}. Under $H_0: \alpha_{i,1} = \cdots = \alpha_{i,q} = 0$,
\begin{equation}
    ARCH_i(q) \overset{d}{\to} \chi^2(q).
\end{equation}
Combine across $i$ via Fisher's method:
\begin{equation}
    \mathcal{F}_{ARCH} = -2 \sum_{i=1}^{24} \log p_i^{ARCH}
    \overset{d}{\to} \chi^2(48).
\end{equation}

\subsection{Multivariate ARCH-LM}

The multivariate extension proposed by Ling and Li (1997) tests whether the matrix
sequence $\{z_t z_t^\top\}$ is serially uncorrelated. Stack the unique elements of
$z_t z_t^\top$ into a vector $\mathrm{vech}(z_t z_t^\top) \in \mathbb{R}^{N(N+1)/2}$,
which has dimension $300$ for $N = 24$. The multivariate ARCH-LM statistic of lag order
$q$ is
\begin{equation}
    MARCH(q) = T \cdot \mathrm{tr}\!\left(\hat{R}^\top \hat{R}\right),
\end{equation}
where $\hat{R}$ is the sample cross-correlation matrix of $\mathrm{vech}(z_t z_t^\top)$
at lags $1$ through $q$. Under $H_0$ (no multivariate ARCH), this statistic is
asymptotically $\chi^2$ with $q[N(N+1)/2]^2$ degrees of freedom. In practice, for
$N = 24$ and $q > 1$ this becomes unwieldy; it is therefore advisable to apply the
univariate ARCH-LM series by series and combine via Fisher's method, reserving the
multivariate version for $N$ small.

%------------------------------------------------------------------
\section{A Complete Testing Protocol}
%------------------------------------------------------------------

We summarize the recommended sequence of tests. Each test addresses a distinct
aspect of the i.i.d.\ null hypothesis for the panel $\{z_{i,t}\}_{i=1,\ldots,24;\;t=1,\ldots,T}$.

\paragraph{Step 1: Marginal distribution.}
Compute $z_{i,t}$ via \eqref{eq:zit}. Apply the Jarque-Bera test to each series and
combine via \eqref{eq:fisher_lb}. Reject if $\mathcal{F}_{JB} > \chi^2_{48}(1-\alpha)$.
Also inspect Q-Q plots of the pooled $z_{i,t}$ against $\mathcal{N}(0,1)$.

\paragraph{Step 2: Serial dependence in levels.}
Apply the Ljung-Box test \eqref{eq:lb} to each $z_{i,t}$ series at lag $L$. Combine
via Fisher \eqref{eq:fisher_lb} to obtain $\mathcal{F}_{LB}$. Reject if
$\mathcal{F}_{LB} > \chi^2_{48}(1-\alpha)$.

\paragraph{Step 3: Serial dependence in squares.}
Repeat Step 2 with $z_{i,t}^2$ in place of $z_{i,t}$ to test for ARCH effects.
Alternatively apply the Engle LM test and combine via $\mathcal{F}_{ARCH}$.

\paragraph{Step 4: Cross-sectional independence.}
Apply the Pesaran CD test \eqref{eq:cd}. Reject if $|CD| > z_{1-\alpha/2}$.
Complement with the full covariance test \eqref{eq:covtest}.

\paragraph{Step 5: Joint serial dependence across all series.}
Apply the multivariate Ljung-Box statistic \eqref{eq:hosking} on the fully
standardized residuals $\tilde{z}_t = \hat{\Gamma}(0)^{-1/2} z_t$ at lag $L$.
Reject if $Q_M(L) > \chi^2_{N^2 L}(1-\alpha)$.

\paragraph{Step 6: Mahalanobis distance analysis.}
Plot $d_t^2$ over time and compare its empirical distribution with $\chi^2(24)$.
Large isolated spikes indicate outliers; systematic patterns indicate model
misspecification. Test i.i.d.-ness of $\{d_t^2\}$ via a univariate Ljung-Box on
the scalar series.

\paragraph{Interpretation.}
\begin{itemize}
    \item Failure at Step 1 alone: the distributional family $h$ is misspecified.
        Switch from Gaussian to Student's $t$ or add skewness.
    \item Failure at Steps 2--3 but not 4: there is residual time-series structure
        in individual series. Consider adding lags to the state equation or allowing
        time-varying $H_t$.
    \item Failure at Step 4 but not 2--3: the covariance adjustment $M C M^\top$
        has not fully captured cross-sectional dependence. Consider increasing $p$
        or allowing a richer structure for $C$.
    \item Failure at Step 5 only: there are cross-serial dependencies not captured
        by individual tests; the factor structure may be too rigid.
    \item Failure at Step 6 with time-clustering: the model fails during turbulent
        periods, consistent with the COVID findings in the paper.
\end{itemize}

%------------------------------------------------------------------
\section{Notes on Estimation Uncertainty}
%------------------------------------------------------------------

All tests above treat $\hat{\Sigma}$, $\hat{\beta}_t$, and $\hat{\nu}$ as if they were
known. Estimation uncertainty inflates the true size of the tests, particularly for small
$T$. Two corrections are available in practice.

\paragraph{Parametric bootstrap.} Simulate $B = 999$ paths
$\{y_t^{(b)}\}_{t=1}^T$ from the estimated model, re-estimate $\psi$ on each, compute
the test statistic, and use the empirical quantile of the simulated statistics as the
critical value. This is the most reliable approach but is computationally expensive.

\paragraph{Degrees-of-freedom correction.} For the Ljung-Box statistic, reduce
the degrees of freedom from $L$ to $L - p_\psi$, where $p_\psi$ is the number of
estimated parameters governing the dynamics. For the Hosking statistic, the correction
is $N^2 L - p_\psi$. These corrections are conservative and may not fully account for
the complexity of the score-driven filter.

%------------------------------------------------------------------
\section{Connection to the Implied Volatility Application}
%------------------------------------------------------------------

In the empirical application of Zou, Lin, and Lucas (2025), the 24 buckets consist of
four time-to-maturity groups (7--45, 45--90, 90--180, 180--360 days) crossed with six
moneyness groups (DOTM put, OTM put, ATM put, ATM call, OTM call, DOTM call). The
representative contract for each bucket is the one closest to the bucket midpoint.

Under the five-factor specification
\begin{equation}
    \log IV_{i,t} = \beta_{1,t} + \beta_{2,t} m_{i,t} + \beta_{3,t} m_{i,t}^2
    + \beta_{4,t} \tau_{i,t} + \beta_{5,t} m_{i,t} \tau_{i,t} + \varepsilon_{i,t},
\end{equation}
the factor loading matrix $M \in \mathbb{R}^{24 \times 5}$ is the design matrix evaluated
at the 24 bucket representative contracts. The adjusted covariance is then
$\hat{\Sigma} = \hat{H} + M\hat{C}M^\top$ with $\hat{H} = \hat{\sigma}^2_\varepsilon I_{24}$
and $\hat{C} \in \mathbb{R}^{5 \times 5}$.

Several specific failure modes are worth checking:

\begin{enumerate}
    \item \textbf{Volatility smile residuals.} Group the 24 $z_{i,t}$ series by maturity
        bucket and check whether the cross-maturity pattern of autocorrelations is uniform.
        Systematic differences across maturity groups signal that the slope factor
        $\beta_{4,t}$ is not flexible enough.
    \item \textbf{COVID period.} Restrict the Ljung-Box and Hosking tests to the COVID
        sub-sample (2020--2022) and compare with the pre-COVID period. The Student's $t$
        model is known to adapt slowly to level shifts; this should be reflected in
        elevated $d_t^2$ values during March--April 2020.
    \item \textbf{Deep out-of-the-money puts.} DOTM put buckets (small $|\Delta|$) are
        the most illiquid and are most likely to exhibit heavy tails or
        intermittent outliers. A series-specific KS test is recommended for these buckets.
\end{enumerate}

%------------------------------------------------------------------
\end{document}